\setcounter{chapter}{4}
\renewcommand{\thetheorem}{\thechapter.\arabic{theorem}}
\setcounter{theorem}{0}

\chapter{Systems of Linear Equations}

\subsection{Fields (Körper)}

Before we solve systems of linear equations, we must define the algebraic structures over which these equations are built. Throughout this course, we will work over a base \qt{field} (in German: \qt{Körper}), usually denoted by $K$. 

A field is a set equipped with two operations (addition $+$ and multiplication $\cdot$) that satisfy a standard set of axioms, including commutativity, associativity, distributivity, and the existence of identities ($0$ and $1$) and inverses.

\begin{example}[Common Fields]
\begin{itemize}
    \item $\mathbb{Q} = \left\{\frac{m}{n} \mid m, n \in \mathbb{Z}, n \ne 0\right\}$ (The rational numbers)
    \item $\mathbb{R}$ (The real numbers)
    \item $\mathbb{C} = \{a + i \cdot b \mid a, b \in \mathbb{R}, i^2 = -1\}$ (The complex numbers)
\end{itemize}
\end{example}

\begin{example}[Finite Fields]
Fields do not have to be infinite. We will frequently use finite fields, where arithmetic is performed modulo a prime number. 
\begin{enumerate}
    \item $\mathbb{F}_2 = \{0, 1\}$. Operations are done modulo $2$:
    \[
    \begin{array}{c|cc}
    + & 0 & 1 \\
    \hline
    0 & 0 & 1 \\
    1 & 1 & 0
    \end{array}
    \qquad\qquad
    \begin{array}{c|cc}
    \cdot & 0 & 1 \\
    \hline
    0 & 0 & 0 \\
    1 & 0 & 1
    \end{array}
    \]
    \item $\mathbb{F}_3 = \{0, 1, 2\}$. Operations are done modulo $3$:
    \[
    \begin{array}{c|ccc}
    + & 0 & 1 & 2 \\
    \hline
    0 & 0 & 1 & 2 \\
    1 & 1 & 2 & 0 \\
    2 & 2 & 0 & 1
    \end{array}
    \qquad\qquad
    \begin{array}{c|ccc}
    \cdot & 0 & 1 & 2 \\
    \hline
    0 & 0 & 0 & 0 \\
    1 & 0 & 1 & 2 \\
    2 & 0 & 2 & 1
    \end{array}
    \]
\end{enumerate}
\end{example}

\subsection{Motivating Example}

Let us solve a familiar system of linear equations over $\mathbb{R}$ using elimination:
\begin{equation*}
\begin{cases}
E_1: \quad x + 3y + 4z = -1 \\
E_2: \quad 2x - y + z = 5 \\
E_3: \quad 3x + 2y + z = 0
\end{cases}
\end{equation*}

We perform operations on the equations to systematically eliminate variables:
\begin{align*}
-2 \cdot E_1 + E_2 \to E_2 \\
-3 \cdot E_1 + E_3 \to E_3 
\end{align*}
This yields an equivalent, simpler system:
\begin{equation*}
\begin{cases}
x + 3y + 4z = -1 \\
-7y - 7z = 7 \\
-7y - 11z = 3
\end{cases}
\end{equation*}

Next, we scale the second equation ($-\frac{1}{7} \cdot E_2 \to E_2$) to get:
\begin{equation*}
\begin{cases}
x + 3y + 4z = -1 \\
y + z = -1 \\
-7y - 11z = 3
\end{cases}
\end{equation*}

Finally, we eliminate $y$ from the third equation ($7 \cdot E_2 + E_3 \to E_3$):
\begin{equation*}
\begin{cases}
x + 3y + 4z = -1 \\
y + z = -1 \\
-4z = -4 \implies z = 1
\end{cases}
\end{equation*}

By backward substitution:
\begin{itemize}
    \item $y + 1 = -1 \implies y = -2$
    \item $x + 3(-2) + 4(1) = -1 \implies x - 2 = -1 \implies x = 1$
\end{itemize}
Thus, the unique solution is $(x,y,z) = (1, -2, 1)$. We can readily verify this by plugging the values back into the original system.

\subsection{General Systems of Linear Equations}

\begin{definition}[Linear Equation]
A \qt{linear equation} in $n$ variables $x_1, \dots, x_n$ over a field $K$ is an equation of the form:
\[
a_1 x_1 + \dots + a_n x_n = b
\]
where $a_1, \dots, a_n \in K$ are the \qt{coefficients} and $b \in K$ is the \qt{constant term}.
\begin{itemize}
    \item If $b = 0$, the equation is called \qt{homogeneous}.
    \item If $b \ne 0$, the equation is called \qt{inhomogeneous}.
\end{itemize}
\end{definition}

\begin{definition}[System of Linear Equations]
A general system of $m$ linear equations in $n$ variables over a field $K$ is written as:
\begin{equation} \label{eq:general_system}
\begin{cases}
A_{1,1}x_1 + A_{1,2}x_2 + \dots + A_{1,n}x_n = b_1 \\
A_{2,1}x_1 + A_{2,2}x_2 + \dots + A_{2,n}x_n = b_2 \\
\vdots \\
A_{m,1}x_1 + A_{m,2}x_2 + \dots + A_{m,n}x_n = b_m
\end{cases}
\end{equation}
\end{definition}

\begin{definition}[Solution]
A \qt{solution} to the system is a sequence of scalars $s_1, \dots, s_n \in K$ that simultaneously satisfies all $m$ equations when we substitute $x_j = s_j$.
\end{definition}

\begin{definition}[Equivalence]
Two systems of linear equations (in the same variables) are called \qt{equivalent} if they have exactly the same set of solutions.
\end{definition}

\subsection{Matrices}

To study systems of linear equations systematically, we strip away the variables and isolate the coefficients into a rectangular array.

\begin{definition}[Matrix]
An $m \times n$ \qt{matrix} over $K$ is a rectangular array consisting of $m \cdot n$ elements of $K$, arranged in $m$ rows and $n$ columns.
\end{definition}

\begin{notation}
We denote a matrix $A$ by its entries:
\[
A = (A_{i,j})_{1 \le i \le m, 1 \le j \le n} \quad \text{or simply} \quad (a_{ij})
\]
Here, $i$ denotes the row index and $j$ denotes the column index.
\end{notation}

We define the set of all such matrices as:
\begin{equation*}
    \M_{m \times n}(K) := \text{the set of all } m \times n \text{ matrices over } K
\end{equation*}

\begin{definition}[Row and Column Vectors]
Special cases of matrices occur when one of the dimensions is $1$:
\begin{enumerate}
    \item \textbf{Column Vectors:} An $n \times 1$ matrix is called a \qt{column vector}. The set of all column vectors of length $n$ is denoted by $K^n := \M_{n \times 1}(K)$.
    \[
    \begin{pmatrix} x_1 \\ x_2 \\ \vdots \\ x_n \end{pmatrix} \in K^n
    \]
    \item \textbf{Row Vectors:} A $1 \times n$ matrix is called a \qt{row vector}. The set of all row vectors of length $n$ is denoted by $K_{\text{row}}^n := \M_{1 \times n}(K)$.
    \[
    (x_1, x_2, \dots, x_n) \in K_{\text{row}}^n
    \]
\end{enumerate}
\end{definition}

\subsection{Matrix-Vector Multiplication}

\begin{definition}[Matrix-Vector Multiplication]
Let $A \in \M_{m \times n}(K)$ and $x \in K^n$. We define the product $A \cdot x$ as a new column vector in $K^m$, computed by taking the dot product of the rows of $A$ with the vector $x$:
\[
\begin{pmatrix} 
A_{1,1} & A_{1,2} & \dots & A_{1,n} \\
A_{2,1} & A_{2,2} & \dots & A_{2,n} \\
\vdots & \vdots & \ddots & \vdots \\
A_{m,1} & A_{m,2} & \dots & A_{m,n} 
\end{pmatrix}
\begin{pmatrix} x_1 \\ x_2 \\ \vdots \\ x_n \end{pmatrix}
=
\begin{pmatrix} 
A_{1,1}x_1 + \dots + A_{1,n}x_n \\
A_{2,1}x_1 + \dots + A_{2,n}x_n \\
\vdots \\
A_{m,1}x_1 + \dots + A_{m,n}x_n 
\end{pmatrix} \in K^m
\]
\end{definition}

Using this powerful notation, the massive system of linear equations \eqref{eq:general_system} can be written elegantly and compactly as a single matrix equation:
\begin{equation} \label{eq:matrix_equation}
A x = b
\end{equation}
where $A \in \M_{m \times n}(K)$ is the \qt{coefficient matrix}, $x \in K^n$ is the vector of variables, and $b \in K^m$ is the \qt{constant vector}.

\subsection{Elementary Row Operations (EROs)}

To solve $A x = b$, we define the \qt{augmented matrix} $(A|b) \in \M_{m \times (n+1)}(K)$ by appending $b$ as an extra column to $A$. We can manipulate this matrix using three basic operations that correspond exactly to operations on the equations themselves.

\begin{definition}[Elementary Row Operations]
Let $E_i$ denote the $i$-th row of the augmented matrix. The three Elementary Row Operations (EROs) are defined and denoted as follows:
Let $E_i$ denote the $i$-th row of the augmented matrix. The three \qt{Elementary Row Operations} (EROs) are defined and denoted as follows:
\begin{enumerate}
    \item \textbf{Type 1:} Multiply row $i$ by a non-zero scalar $\lambda \in K \setminus \{0\}$. \\
    Notation: $\lambda \cdot E_i \to E_i$
    \item \textbf{Type 2:} Add to row $j$ a scalar multiple of row $i$. \\
    Notation: $\lambda \cdot E_i + E_j \to E_j$
    \item \textbf{Type 3:} Swap row $i$ and row $j$. \\
    Notation: $E_i \leftrightarrow E_j$
\end{enumerate}
\end{definition}

\begin{theorem}
Performing Elementary Row Operations on the augmented matrix $(A|b)$ does NOT change the set of solutions to the corresponding system of linear equations. In other words, if $(A|b)$ is transformed into $(A'|b')$ via EROs, then $Ax = b$ and $A'x = b'$ are equivalent systems.
\end{theorem}

\begin{proof}[Idea of the Proof]
Each of the three EROs is strictly reversible. If you apply an operation, you can apply its inverse to get the original system back. Therefore, no solutions are lost, and no false solutions are created.
\end{proof}

\subsection{Further Computational Examples}

\begin{example}[A system with infinitely many solutions]
Consider the following system:
\begin{equation*}
\begin{cases}
x_1 - 2x_2 + x_3 + x_4 = 1 \\
2x_1 - 4x_2 + x_3 + 3x_4 = 2 \\
x_1 - 2x_2 + 2x_3 - x_4 = 1
\end{cases}
\end{equation*}
We write the augmented matrix and apply EROs:
\[
\begin{pmatrix}
1 & -2 & 1 & 1 & \big| & 1 \\
2 & -4 & 1 & 3 & \big| & 2 \\
1 & -2 & 2 & -1 & \big| & 1 
\end{pmatrix}
\xrightarrow{\begin{subarray}{l} -2 \cdot E_1 + E_2 \to E_2 \\ -1 \cdot E_1 + E_3 \to E_3 \end{subarray}}
\begin{pmatrix}
1 & -2 & 1 & 1 & \big| & 1 \\
0 & 0 & -1 & 1 & \big| & 0 \\
0 & 0 & 1 & -2 & \big| & 0 
\end{pmatrix}
\]
\[
\xrightarrow{1 \cdot E_2 + E_3 \to E_3}
\begin{pmatrix}
1 & -2 & 1 & 1 & \big| & 1 \\
0 & 0 & -1 & 1 & \big| & 0 \\
0 & 0 & 0 & -1 & \big| & 0 
\end{pmatrix}
\]
From the third row, $-x_4 = 0 \implies x_4 = 0$. \\
From the second row, $-x_3 + x_4 = 0 \implies x_3 = 0$. \\
Substituting into the first row gives $x_1 - 2x_2 + 0 + 0 = 1 \implies x_1 = 1 + 2x_2$.

We have a free variable. Let $x_2 = \alpha \in \mathbb{R}$. The general solution is:
\[
x = \begin{pmatrix} x_1 \\ x_2 \\ x_3 \\ x_4 \end{pmatrix} = \begin{pmatrix} 1 + 2\alpha \\ \alpha \\ 0 \\ 0 \end{pmatrix} = \begin{pmatrix} 1 \\ 0 \\ 0 \\ 0 \end{pmatrix} + \alpha \begin{pmatrix} 2 \\ 1 \\ 0 \\ 0 \end{pmatrix}
\]
\end{example}

\begin{example}[Conditions for consistency]
Let's see how solutions behave for a general choice of constants $b_1, b_2, b_3 \in \mathbb{R}$:
\[
\begin{pmatrix}
1 & -2 & 1 \\
2 & 1 & 1 \\
0 & 5 & -1 
\end{pmatrix}
\begin{pmatrix} x_1 \\ x_2 \\ x_3 \end{pmatrix}
= \begin{pmatrix} b_1 \\ b_2 \\ b_3 \end{pmatrix}
\]
We construct the augmented matrix and reduce:
\[
\begin{pmatrix}
1 & -2 & 1 & \big| & b_1 \\
2 & 1 & 1 & \big| & b_2 \\
0 & 5 & -1 & \big| & b_3 
\end{pmatrix}
\xrightarrow{-2 \cdot E_1 + E_2 \to E_2}
\begin{pmatrix}
1 & -2 & 1 & \big| & b_1 \\
0 & 5 & -1 & \big| & b_2 - 2b_1 \\
0 & 5 & -1 & \big| & b_3 
\end{pmatrix}
\]
\[
\xrightarrow{-1 \cdot E_2 + E_3 \to E_3}
\begin{pmatrix}
1 & -2 & 1 & \big| & b_1 \\
0 & 5 & -1 & \big| & b_2 - 2b_1 \\
0 & 0 & 0 & \big| & b_3 - b_2 + 2b_1 
\end{pmatrix}
\]
For this system to have \textit{any} solutions, the bottom row cannot result in a contradiction ($0 = \text{non-zero}$). Therefore, the system is consistent if and only if:
\[
b_3 - b_2 + 2b_1 = 0
\]
\end{example}